\documentclass[border=10pt]{standalone}
\usepackage{tikz}
\usetikzlibrary{chains, calc}
\begin{document}
% Top-frame bits: real/visual. Visual = real except for the
% payload, which is compressed and gets a dashed break.
\def\PreBits{89}\def\HdrBits{32}
\def\PldBits{12000}\def\PldVisualBits{100}\def\PCRCBits{16}
\def\BreakFrac{0.40}    % fraction of payload visible-width that's the gap
% Header detail (bits each)
\def\PlBits{11}\def\SBits{2}\def\DBits{2}
\def\TBits{2}\def\MBits{3}\def\QBits{7}\def\GBits{3}\def\CBits{8}
% Drawing dimensions (cm)
\def\FrameWidth{17}\def\BoxHeight{0.7}\def\ZoomGap{0.8}

\def\bitblock#1#2#3{%
  \pgfmathsetmacro{\W}{#1*\bitscale}%
  \pgfmathtruncatemacro{\firstbit}{\bitsofar}%
  \pgfmathtruncatemacro{\lastbit}{\bitsofar+#1-1}%
  \node[box, on chain, minimum width=\W cm] (#3) {#2};%
  \node[bit, anchor=south west] at (#3.north west) {\firstbit};%
  \ifnum#1>1\node[bit, anchor=south east] at (#3.north east) {\lastbit};\fi
  \xdef\bitsofar{\the\numexpr\bitsofar+#1\relax}%
}
\def\bitblockbreak#1#2#3#4#5{%
  \pgfmathsetmacro{\W}{#2*\bitscale}%
  \pgfmathtruncatemacro{\firstbit}{\bitsofar}%
  \pgfmathtruncatemacro{\lastbit}{\bitsofar+#1-1}%
  \pgfmathsetmacro{\gapstart}{(1-#5)/2}%
  \pgfmathsetmacro{\gapend}{1-\gapstart}%
  \node[box, on chain, draw=none, minimum width=\W cm] (#4) {#3};%
  \path
    ($(#4.south west)!\gapstart!(#4.south east)$) coordinate (#4-gap-bl)
    ($(#4.south west)!\gapend!(#4.south east)$) coordinate (#4-gap-br)
    ($(#4.north west)!\gapstart!(#4.north east)$) coordinate (#4-gap-tl)
    ($(#4.north west)!\gapend!(#4.north east)$) coordinate (#4-gap-tr);
  \draw[solid]
    (#4-gap-bl) -- (#4.south west) -- (#4.north west) -- (#4-gap-tl)
    (#4-gap-br) -- (#4.south east) -- (#4.north east) -- (#4-gap-tr);
  \draw[densely dashed] (#4-gap-tl) -- (#4-gap-tr)  (#4-gap-bl) -- (#4-gap-br);
  \node[bit, anchor=south west] at (#4.north west) {\firstbit};%
  \node[bit, anchor=south east] at (#4.north east) {\lastbit};%
  \xdef\bitsofar{\the\numexpr\bitsofar+#1\relax}%
}

\begin{tikzpicture}[
  box/.style  = {draw, inner sep=0pt, outer sep=0pt,
                 minimum height=\BoxHeight cm,
                 font=\footnotesize, align=center,
                 execute at begin node=\strut},
  bit/.style  = {font=\scriptsize},
  zoom/.style = {densely dashed, gray},
]
  \pgfmathsetmacro{\TopK}{\FrameWidth/(\PreBits+\HdrBits+\PldVisualBits+\PCRCBits)}
  \pgfmathsetmacro{\BotK}{\FrameWidth/(\PlBits+\SBits+\DBits+\TBits+\MBits+\QBits+\GBits+\CBits)}
  \pgfmathsetmacro{\BotY}{-\ZoomGap-\BoxHeight}

  % --- Top frame: Preamble, Header, Payload (with break) ----------
  \begin{scope}[start chain=top, node distance=0pt]
    \let\bitscale\TopK \xdef\bitsofar{0}
    \node[on chain, inner sep=0pt, outer sep=0pt, minimum size=0pt,
          anchor=west] at (0,0) {};
    \bitblock{\PreBits}{Preamble}{pre}
    \bitblock{\HdrBits}{Header}{hdr}
    \bitblockbreak{\PldBits}{\PldVisualBits}{Payload}{pld}{\BreakFrac}
    \pgfmathsetmacro{\PCRCW}{\PCRCBits*\bitscale}
    \node[box, on chain, minimum width=\PCRCW cm] (pcrc) {CRC};
    \node[bit, anchor=south] at (pcrc.north) {+\PCRCBits};
  \end{scope}

  % --- Bottom frame: header detail --------------------------------
  \begin{scope}[start chain=bot, node distance=0pt,
                yshift=\BotY cm]
    \let\bitscale\BotK \xdef\bitsofar{0}
    \node[on chain, inner sep=0pt, outer sep=0pt, minimum size=0pt,
          anchor=west] at (0,0) {};
    \foreach \bits/\lbl/\id in {%
      \PlBits/Payload length/pl,
      \SBits/Src/s,
      \DBits/Dest/d,
      \TBits/Type/t,
      \MBits/Mod./m,
      \QBits/Sequence Number/q,
      \GBits/Coding/g,
      \CBits/CRC/c%
    }{%
      \bitblock{\bits}{\lbl}{\id}%
    }
  \end{scope}

  % --- Zoom-out lines: Header (top) -> bottom frame ---------------
  \draw[zoom] (hdr.south west) -- (pl.north west);
  \draw[zoom] (hdr.south east) -- (c.north east);
\end{tikzpicture}
\end{document}
